\section{Preliminaries}\label{sec:prelim}

We recall the relevant notions from propositional proof complexity and
chain-of-thought reasoning, and fix notation used throughout the paper.

\subsection{Propositional proof complexity}\label{sec:prelim-proof}

A \emph{literal} is a Boolean variable $x$ or its negation $\bar{x}$.  A
\emph{clause} is a disjunction of literals.  A \emph{conjunctive normal form}
(CNF) formula $F$ is a conjunction of clauses.  We write $\operatorname{Var}(F)$
for the set of variables appearing in~$F$ and set $n = \abs{\operatorname{Var}(F)}$.

\begin{definition}[Resolution proof system]\label{def:resolution}
The \emph{resolution} rule derives a new clause from two parent clauses sharing
a complementary literal:
\[
\frac{C \vee x \qquad D \vee \bar{x}}{C \vee D}\,.
\]
A \emph{resolution refutation} (or \emph{proof}) of an unsatisfiable CNF
formula~$F$ is a sequence of clauses $C_1, C_2, \ldots, C_S$ such that each
$C_i$ is either a clause of~$F$ or is derived from two earlier clauses by the
resolution rule, and $C_S$ is the empty clause~$\bot$.
\end{definition}

\begin{definition}[Proof complexity measures]\label{def:complexity-measures}
Let $\pi = (C_1, \ldots, C_S)$ be a resolution refutation of~$F$.
\begin{enumerate}[label=(\roman*)]
\item The \emph{length} (or \emph{size}) of $\pi$ is $S(\pi) = S$.
\item The \emph{width} of~$\pi$ is
  $w(\pi) = \max_{1 \le i \le S} \abs{C_i}$, where $\abs{C_i}$ denotes the
  number of literals in~$C_i$.
\item The \emph{depth} of~$\pi$ is the height of its derivation DAG.
\item The \emph{space} of~$\pi$ is the maximum number of clauses that must be
  simultaneously retained in memory when executing the proof in a
  sequential blackboard model.
\end{enumerate}
We write $S(F)$, $w(F \vdash \bot)$, $d(F)$, and $\operatorname{Sp}(F)$ for
the respective minima over all refutations of~$F$.
\end{definition}

The following classical results relate these measures.

\begin{theorem}[Ben-Sasson and Wigderson~{\cite{bensasson2001short}}]
\label{thm:bsw}
For any unsatisfiable CNF formula~$F$ on $n$~variables with initial clause
width $w(F)$,
\[
w(F \vdash \bot) \le w(F) + O\!\left(\sqrt{n \ln S(F)}\right).
\]
Equivalently, $S(F) \ge \exp\!\left(\Omega\!\left(\frac{(w(F \vdash \bot) -
w(F))^2}{n}\right)\right)$.
\end{theorem}

We use the following well-known formula families as test instances.

\begin{definition}[Tseitin formulas]\label{def:tseitin}
Let $G = (V, E)$ be a connected graph and $\sigma \colon V \to \{0,1\}$ a
labeling with $\sum_{v \in V} \sigma(v) \equiv 1 \pmod{2}$.  The
\emph{Tseitin formula} $\operatorname{Ts}(G, \sigma)$ introduces a variable
$x_e$ for each edge $e \in E$ and, for each vertex $v$, a set of clauses
encoding $\bigoplus_{e \ni v} x_e = \sigma(v)$.  These formulas are always
unsatisfiable because the total parity constraint is odd.
\end{definition}

\begin{definition}[Pigeonhole formulas]\label{def:php}
The \emph{pigeonhole principle} formula $\operatorname{PHP}^m_n$ (with $m > n$)
has variables $p_{i,j}$ for $1 \le i \le m$, $1 \le j \le n$, with clauses
encoding ``each pigeon maps to at least one hole'' and ``no two pigeons share a
hole.''  Haken~\cite{haken1985intractability} proved that
$S(\operatorname{PHP}^{n+1}_n) = 2^{\Omega(n)}$.
\end{definition}

\begin{definition}[Pebbling formulas]\label{def:pebbling}
Let $G = (V, E)$ be a directed acyclic graph with a unique sink~$t$.  The
\emph{pebbling formula} $\operatorname{Peb}(G)$ has a variable $x_v$ for each
$v \in V$, with clauses encoding: (i)~source vertices are true; (ii)~if all
predecessors of $v$ are true then $v$ is true; and (iii)~the sink is false.
The space complexity of $\operatorname{Peb}(G)$ equals the black pebbling
number of~$G$, plus a constant~\cite{bensasson2001short}.
\end{definition}

\subsection{Chain-of-thought verification}\label{sec:prelim-cot}

We model chain-of-thought verification as a sequential stochastic process.  An
LLM verifier receives a proof and processes it over discrete \emph{reasoning
steps} $t = 1, 2, \ldots$  At each step, the verifier updates its assessment
of the proof's correctness.

\begin{definition}[Process reward model]\label{def:prm}
A \emph{process reward model} (PRM) assigns a correctness probability
$r_t \in [0,1]$ to each step~$t$ of a solution or proof trace.  The
solution-level score is $\prod_{t=1}^T r_t$, where $T$ is the total number of
steps~\cite{lightman2023lets}.
\end{definition}

\begin{definition}[Self-consistency]\label{def:self-consistency}
Given $N$ independent CoT verification traces for a proof~$F$, each
producing a binary judgment $Y_i \in \{0,1\}$ with
$\Pr[Y_i = 1] = p$, the \emph{self-consistency} (majority vote) decision is
\[
\hat{Y}_N = \mathbf{1}\!\left[\sum_{i=1}^N Y_i > N/2\right].
\]
When $p > 1/2$, the majority vote accuracy exceeds~$p$ and converges to~$1$ as
$N \to \infty$~\cite{wang2022self}.
\end{definition}

\subsection{Notation}\label{sec:notation}

Throughout the paper, $F$ denotes an unsatisfiable CNF formula, $S(F)$ its
minimum resolution proof length, $w(F \vdash \bot)$ its minimum proof width,
$d(F)$ its minimum proof depth, and $\operatorname{Sp}(F)$ its minimum proof
space.  We write $\ln$ for the natural logarithm and $\log$ for the base-$2$
logarithm.  The notation $f = O(g)$, $f = \Omega(g)$, and $f = \Theta(g)$
carries the standard meaning.  The Spearman rank correlation coefficient is
denoted~$\rho$.  All probabilities are with respect to the randomness in the
CoT generation process.
